% Jawaban publik bernomor ganjil untuk Bab 3 Bagian 3.4 (latihan 29, 31, 33, 35).
% Diambil hanya dari lapisan eoceSolutions publik yang dipatok; tidak ada jawaban terbatas yang ditambahkan.

% 29

\eocesol{(a)~13.
(b)~Tidak, 27 mahasiswa ini bukan sampel acak dari populasi mahasiswa 
universitas tersebut. Sebagai contoh, dapat dikatakan bahwa proporsi perokok di antara 
mahasiswa yang pergi ke pusat kebugaran pada pukul 9 pagi pada hari Sabtu lebih rendah daripada 
proporsi perokok di universitas secara keseluruhan.}

% 31

\eocesol{(a)~E(X) = 3.59. SD(X) = 9.64.
(b)~E(X) = -1.41. SD(X) = 9.64.
(c)~Tidak, laba bersih yang diharapkan bernilai negatif, sehingga secara rata-rata Anda diperkirakan akan kehilangan uang.}

% 33

\eocesol{Nilainya meningkat 5\%.}

% 35

\eocesol{E = -0.0526. SD = 0.9986.}
